\section{Bosonic chiral differential operators}
\label{sec:chiral-operators}

A local holomorphic coordinate system on $Y$ has a free
beta--gamma vertex algebra. Coordinate changes require quantum
correction terms, and the corrected maps can fail to compose
on triple overlaps. The resulting obstruction lies in closed
holomorphic cohomology. The curve periods of
Section~\ref{sec:tangent-curves} make this obstruction nonzero.
A separate obstruction prevents the standard local Virasoro
fields from gluing.

We retain the same selected threefold and fixed gluing data.
The vertex variable below is a formal one-dimensional coordinate.
It is distinct from the three target coordinates on $Y$.

\subsection{The analytic gerbe and its band}

On a holomorphic coordinate chart $U$ with coordinates
$b=(b^1,b^2,b^3)$, let $\mathcal D^{\mathrm{ch}}_{U,b}$ be the
analytic free bosonic vertex algebra with generating fields
$b^i(\xi)$ of weight zero and $a^i(\xi)$ of weight one.
We fix the operator product expansion
\begin{equation}\label{ch:eq:generator-OPE}
       a^i(\xi)b^j(\eta)\sim\frac{\delta_{ij}}{\xi-\eta},
\end{equation}
with the other generator products regular.
The symbol $\sim$ records the singular part.
Holomorphic functions of the $b^i$ form the weight-zero sheaf
$\mathcal O_U$, and its constant $1$ is the vacuum.
We write $\mathsf T$ for vertex translation and $d_Y$ for
holomorphic exterior differentiation.
The analytic construction and holomorphic functional calculus
are those of \cite[Section~2]{GMSGerbes}.

The weight-one filtration has one-form submodule and tangent
quotient:
\[
       0\longrightarrow\Omega_U^1
       \longrightarrow\mathcal D^{\mathrm{ch}}_{U,b,1}
       \longrightarrow T_U\longrightarrow0.
\]
The one-form $db^i$ corresponds to $\mathsf T b^i$.
For a coordinate change $\widetilde b=g(b)$, use the Jacobian
convention $J=(\partial_i g^j)$.
Its tangent term uses $J^{-t}$, as in
\cite[equations~(4.4)--(4.5)]{GMSGerbes}.
An admissible lift must preserve both the mixed OPE
\eqref{ch:eq:generator-OPE} and the regular self-OPEs.

The analytic groupoid of these local objects is a gerbe
$\mathfrak D_Y^{\mathrm{ch}}$ with band
$\Omega_Y^{2,\mathrm{cl}}$, the sheaf of closed holomorphic
two-forms \cite[Theorem~3 and Section~6]{GMSGerbes}.
We use its published band normalization.
In coordinates, a change of lift over the identity is represented
by a skew matrix $h=(h^{ki})$ and acts by
\begin{equation}\label{ch:eq:skew-twist}
       b^i\longmapsto b^i,\qquad
       a^i\longmapsto a^i+\sum_k h^{ki}(b)\mathsf T b^k.
\end{equation}
The remaining OPE condition says that the associated holomorphic
two-form is closed. This is the kernel identification in
\cite[Section~6]{GMSGerbes}.

For clarity, choose an analytic cover with coordinate transitions
$b_j=g_{ij}(b_i)$.
On sufficiently small overlaps, write
$\phi_{ij}:\mathcal D_j^{\mathrm{ch}}\to\mathcal D_i^{\mathrm{ch}}$
for an admissible lift expressing the $j$-fields in the $i$-fields.
Composition is right to left. The discrepancy
\[
                   c_{ijk}=\phi_{ij}\phi_{jk}\phi_{ik}^{-1}
\]
is a band automorphism of the $i$-object.
These discrepancies give the class
\begin{equation}\label{ch:eq:gerbe-class}
       \omega_{\mathrm{GMS}}
       =c(\mathfrak D_Y^{\mathrm{ch}})
          \in H^2(Y,\Omega_Y^{2,\mathrm{cl}}).
\end{equation}
Changing the pairwise lifts by closed two-forms changes the
cocycle by their \v{C}ech coboundary.
Vanishing of this class is called \emph{neutrality} of the gerbe.
It is the obstruction to existence of a graded bosonic CDO sheaf.
The nonvanishing below will exclude compatible lifts on the full atlas.

\subsection{The Hodge image and the curve criterion}

The same exact sequence \eqref{cq:eq:Poincare} that classifies
Courant extensions also controls the gerbe class, in the next
cohomological degree. Denote its degree-two forgetful map by
\[
 \iota_2:H^2(Y,\Omega_Y^{2,\mathrm{cl}})
                     \longrightarrow H^2(Y,\Omega_Y^2).
\]
It is a map of complex vector spaces. The differential in that
sheaf sequence is not an $\mathcal O_Y$-module map.

\begin{proposition}\label{ch:prop:band-cohomology}
One has
\begin{equation}\label{ch:eq:band-cohomology}
 \begin{gathered}
 H^0(Y,\Omega_Y^{2,\mathrm{cl}})=0,\qquad
 H^1(Y,\Omega_Y^{2,\mathrm{cl}})\simeq\mathbb C,\\
 \iota_2\text{ is injective},\qquad
 \im\iota_2
    =\ker\bigl(d_1:H_{\bar\partial}^{2,2}
                           \to H_{\bar\partial}^{3,2}\bigr)
    =j_{BC}(L).
 \end{gathered}
\end{equation}
In particular, $H^2(Y,\Omega_Y^{2,\mathrm{cl}})$ is a
complex line.
\end{proposition}

\begin{proof}
The Hodge table \eqref{hr:eq:hodge-table} gives
$H^0(\Omega_Y^2)=H^0(\Omega_Y^3)=H^1(\Omega_Y^3)=0$.
The long exact sequence of \eqref{cq:eq:Poincare} gives the
first two assertions and injectivity of $\iota_2$.
Its image is the kernel of the displayed map induced by $d_Y$,
which is the first Fr\"olicher differential.
The row is exact by Proposition~\ref{hr:prop:frolicher}.
Its kernel is the one-dimensional image of
$\partial:H_{\bar\partial}^{1,2}\to H_{\bar\partial}^{2,2}$,
namely $j_{BC}(L)$.
\end{proof}

We next fix the characteristic-class normalization.
Gorbounov--Malikov--Schechtman use raw $d\log$ transition
cocycles. Their Hodge image formula is
\begin{equation}\label{ch:eq:published-image}
       \iota_2(\omega_{\mathrm{GMS}})
                    =2\operatorname{ch}^{\mathrm{Hodge}}_2(T_Y)
\end{equation}
\cite[Claim~3.1 and Theorem~6.2]{GMSGerbes}.
In the normalization \eqref{tc:eq:character-normalization},
\begin{equation}\label{ch:eq:normalization-comparison}
 \operatorname{ch}^{\mathrm{Hodge}}_2(T_Y)
       =(2\pi i)^2j_{BC}\bigl(\operatorname{ch}_{2,BC}(T_Y)\bigr).
\end{equation}
Indeed, each Atiyah factor in the quadratic characteristic
polynomial is divided by $2\pi i$ for the normalized Chern
character. A simultaneous sign change in the degree-one
\v{C}ech--Dolbeault convention cancels in degree two.

Only the Hodge image formula is needed here.
Proposition~\ref{ch:prop:band-cohomology} recovers the refined
class from that image on this particular $Y$.
This does not require a general refined formula beyond the
published Hodge statement.
Nor does it declare the particular matrix in
\cite[equation~(5.3)]{GMSGerbes} to be admissible on every
three-dimensional overlap: that formula alone solves only the
first displayed OPE condition.

\begin{theorem}\label{ch:thm:curve-criterion}
Retain the periods $a_j$, their sum $A$, and the class $w$
from \eqref{tc:eq:period-normalization}--\eqref{tc:eq:curve-periods}.
Then
\begin{equation}\label{ch:eq:curve-image}
 \begin{aligned}
 \iota_2(\omega_{\mathrm{GMS}})
   &=-2(2\pi i)^2j_{BC}\!\left(\sum_j[C_j]_{BC}\right)\\
   &=2(2\pi i)^2 A\,j_{BC}(w).
 \end{aligned}
\end{equation}
Consequently,
\begin{equation}\label{ch:eq:neutrality-criterion}
 \omega_{\mathrm{GMS}}=0
 \quad\Longleftrightarrow\quad c_{2,BC}(T_Y)=0
 \quad\Longleftrightarrow\quad \sum_j[C_j]_{BC}=0
 \quad\Longleftrightarrow\quad A=0.
\end{equation}
If the gerbe is neutral, $Y$ admits no pluriclosed Hermitian
metric. Equivalently, such a metric would imply
$\omega_{\mathrm{GMS}}\ne0$.
\end{theorem}

\begin{proof}
Theorem~\ref{bc:thm:ring} gives $c_{1,BC}(T_Y)^2=0$.
Thus
\[
        \operatorname{ch}_{2,BC}(T_Y)
                           =-c_{2,BC}(T_Y)
                           =-\sum_j[C_j]_{BC}.
\]
Combine this with \eqref{ch:eq:published-image},
\eqref{ch:eq:normalization-comparison} and
\eqref{tc:eq:curve-coordinates}.
Since $j_{BC}(\overline w)=0$, this proves
\eqref{ch:eq:curve-image}.

The map $\iota_2$ is injective, and $j_{BC}$ is injective
on the real Bott--Chern plane.
The current sum is real, so its Hodge image vanishes precisely
when its Bott--Chern class vanishes.
The independent classes $w,\overline w$ then give the equivalence
with $A=0$.
The metric consequence is
Proposition~\ref{tc:prop:global-potential}, applied to the
nonzero positive current $\sum_j[C_j]$.
\end{proof}

\begin{corollary}\label{ch:cor:nonzero-gerbe}
For the selected threefold and the normalization by $\xi$,
\begin{equation}\label{ch:eq:nonzero-gerbe}
 \iota_2(\omega_{\mathrm{GMS}})
              =2(2\pi i)^2j_{BC}(w)\ne0.
\end{equation}
There is no globally defined graded bosonic CDO sheaf on $Y$
with the specified tangent action and the published band.
\end{corollary}

\begin{proof}
Theorem~\ref{tc:thm:evaluated-curves} gives $A=1$ and
$c_{2,BC}(T_Y)=-\rho_{BC}$.
The map $j_{BC}$ kills $\overline w$ and is injective on
$\mathbb Cw$. Equation~\eqref{ch:eq:curve-image} therefore
gives \eqref{ch:eq:nonzero-gerbe}.
The gerbe class is nonzero, so the gerbe is not neutral.
\end{proof}

Rescaling the coefficient band gives
\[
 \widehat\omega=(2\pi i)^{-2}\omega_{\mathrm{GMS}},\qquad
 \iota_2(\widehat\omega)=-2j_{BC}(c_{2,BC}(T_Y))
                                      =2j_{BC}(w)\ne0.
\]
This rescaling changes neither the published local lifts nor the
vanishing criterion.
It keeps the raw gerbe class distinct from the normalized
Bott--Chern characteristic class.

For a neutral gerbe with this band, the isomorphism classes of
objects form a torsor under $H^1(Y,\Omega_Y^{2,\mathrm{cl}})$.
The band automorphism group of an object is
$H^0(Y,\Omega_Y^{2,\mathrm{cl}})$
\cite[Section~3]{GMSGerbes}.
Thus the former is a torsor under a complex line, and the latter
is trivial by \eqref{ch:eq:band-cohomology}.
Neutrality would not choose a distinguished object.
On the selected $Y$, Corollary~\ref{ch:cor:nonzero-gerbe}
excludes neutrality and hence the existence of such objects.

\subsection{The standard Virasoro obstruction}

The standard local Virasoro field is
\[
                        L_i=\sum_{r=1}^3 a_i^r\mathsf T b_i^r,
\]
with the vertex normal ordering of the local free algebra.
For the overlap direction $b_j=g_{ij}(b_i)$, put
$J_{ij}=(\partial_r g_{ij}^s)$.
The published transformation law reads
\begin{equation}\label{ch:eq:Virasoro-overlap}
       \phi_{ij}(L_j)-L_i
                  =-\frac12\mathsf T^2\log\det J_{ij}
\end{equation}
\cite[Theorem~9.1]{GMSGerbes}.
The right side is independent of the local logarithm branch.
We use this law for admissible lifts.
The calculation below proves its one-coordinate instance directly.

For a line bundle $L$ with frame transitions $g_{ij}$, let
$a(L)\in H^1(Y,\Omega_Y^{1,\mathrm{cl}})$ denote its raw
$d_Y\log g_{ij}$ class.
Write $a(T_Y)=a(\det T_Y)$.
Canonical frames transform by
\[
 db_j^1\wedge db_j^2\wedge db_j^3
        =(\det J_{ij})\,db_i^1\wedge db_i^2\wedge db_i^3.
\]
Thus $d_Y\log\det J_{ij}$ represents $a(K_Y)$, and
$a(T_Y)=-a(K_Y)$.
Under the one-form state identification
$\sum_r h_rdb^r\mapsto\sum_r h_r\mathsf T b^r$,
equation~\eqref{ch:eq:Virasoro-overlap} gives the translated
cocycle of
\begin{equation}\label{ch:eq:conformal-class}
                 -\tfrac12a(K_Y)=\tfrac12a(T_Y).
\end{equation}
This fixes its sign in the stated overlap direction.

Define $\mathcal M$ to be the unique line bundle satisfying
\[
                \mathcal M^{\otimes12}=f^*\mathcal O_B(1),
                \qquad \kappa=a(\mathcal M).
\]
Existence and uniqueness follow from
$\operatorname{Pic}(Y)\simeq(\mathbb C,+)$ in
Corollary~\ref{nf:cor:picard}.
This tensor-power normalization does not specify the coordinate
of $f^*\mathcal O_B(1)$ in the coherent parameter $\lambda$
used there.

\begin{theorem}\label{ch:thm:Virasoro-obstruction}
The standard Virasoro obstruction is
\begin{equation}\label{ch:eq:three-kappa}
                     \tfrac12a(T_Y)=3\kappa\ne0.
\end{equation}
For any graded bosonic CDO object that exists on $Y$, its
standard local Virasoro fields cannot be glued, even after
local improvements by translated closed holomorphic one-forms.
Closed two-form changes of admissible lifts do not alter this
obstruction.
\end{theorem}

\begin{proof}
The canonical formula \eqref{nf:eq:canonical} and
$4F_4\sim f^*(1)$ give
\[
             K_Y^{\otimes2}\simeq f^*\mathcal O_B(-1)
                               \simeq\mathcal M^{-12}.
\]
The Picard group has no torsion, so $K_Y\simeq\mathcal M^{-6}$.
Equation~\eqref{ch:eq:conformal-class} now gives $3\kappa$.

The line $f^*\mathcal O_B(1)$ is nontrivial.
Its two base sections pull back to independent sections because
$f$ is surjective. A trivial line on compact connected $Y$
would have only the one-dimensional space of constant sections.
Therefore $\mathcal M$ is also nontrivial.

The holomorphic Poincar\'e lemma gives
\[
       0\longrightarrow\mathbb C\longrightarrow\mathcal O_Y
          \xrightarrow{\ d_Y\ }\Omega_Y^{1,\mathrm{cl}}
          \longrightarrow0.
\]
Since $H^1(Y;\mathbb C)=H^2(Y;\mathbb C)=0$, its cohomology
map
\[
       d_Y:H^1(Y,\mathcal O_Y)
                 \xrightarrow{\ \sim\ }H^1(Y,\Omega_Y^{1,\mathrm{cl}})
\]
is an isomorphism.
For the exponential convention $h\mapsto e^{2\pi i h}$,
the class $a(L)$ is $2\pi i$ times this map on the additive
Picard class of $L$. Nontriviality of $\mathcal M$ therefore
proves $\kappa\ne0$.

An improvement $L_i+\mathsf T\alpha_i$, with $\alpha_i$ a
closed holomorphic one-form state, changes
\eqref{ch:eq:conformal-class} by a coboundary.
Translation is injective on such one-form states:
the coefficients of the independent modes $b^r_{-2}$ in
$\mathsf T(\sum_r h_r(b)b^r_{-1}\mathbf1)$ are $2h_r$.
Thus the nonzero cohomology class prevents gluing even with
these improvements.

Finally, the extra contribution to $L$ from
\eqref{ch:eq:skew-twist} is
\[
                   \sum_{k,r}h^{kr}(b)\mathsf T b^k\mathsf T b^r=0.
\]
The $b$-fields and coefficient functions have regular products,
and the coefficient matrix is skew.
Hence such twists preserve the Virasoro transformation law and
its obstruction.
\end{proof}

Existence of a graded bosonic object and gluing of its standard
conformal fields are therefore distinct questions.
The canonical chiral de Rham construction adds fermionic fields
and gives a sheaf in a larger vertex-superalgebra system
\cite[Section~8]{GMSGerbes}.
That positive construction does not choose a pure bosonic object
or remove the obstruction \eqref{ch:eq:three-kappa} in that object.

\subsection{A quantum coordinate change on an actual chart}

Choose a simply connected logarithmic chart of the nonzero line
coordinate $v_{\mathrm{fib}}$ on the regular part of the quotient
used in \eqref{nf:eq:regular-volume}.
Restrict it so that
\[
                  b=\log v_{\mathrm{fib}},\qquad g(b)=e^b
\]
are inverse coordinate descriptions. The other two local
coordinates are fixed spectators.
All derivatives in the next calculation are with respect to $b$,
and $\xi$ remains the vertex variable.

For the changed bosonic pair, use modes
\begin{equation}\label{ch:eq:modes}
 \begin{gathered}
 a(\xi)=\sum_{m\in\mathbb Z}a_m\xi^{-m-1},\qquad
 b(\xi)=\sum_{n\in\mathbb Z}b_n\xi^{-n},\\
                         [a_m,b_n]=\delta_{m,-n}.
 \end{gathered}
\end{equation}
Functions of $b_0$ act on the vacuum by multiplication.
For any locally invertible one-coordinate change $g$, put
\begin{equation}\label{ch:eq:local-lift}
 \begin{gathered}
 f=\frac1{g'},\qquad
 h=\frac{(g'')^2}{2(g')^3}=\frac{(f')^2}{2f},\\
 B=g(b_0)\mathbf1,\qquad
 A=a_{-1}f(b_0)\mathbf1+b_{-1}h(b_0)\mathbf1.
 \end{gathered}
\end{equation}
The function $f$ in this local formula is unrelated to the
global fibration map.
The order of the two modes and their coefficient functions is
exactly that in \cite[equations~(6.2) and~(9.3)]{GMSGerbes}.
Write $S(\xi)=\mathsf Y(S,\xi)$ for the field of a state $S$.

\begin{proposition}\label{ch:prop:local-OPE}
The states in \eqref{ch:eq:local-lift} satisfy
\[
       A(\xi)B(0)\sim\xi^{-1},\qquad
       A(\xi)A(0)\text{ and }B(\xi)B(0)\text{ are regular}.
\]
Together with the two unchanged spectator pairs, they satisfy
all the free generator OPEs on this chart.
\end{proposition}

\begin{proof}
Write $p=a_{-1}$, $q=b_{-1}$ and $r=b_{-2}$ for the
corresponding creation operators on Fock states.
Negative modes multiply, while positive modes act by
\begin{equation}\label{ch:eq:Fock-action}
       a_0=\frac{\partial}{\partial b_0},\qquad
       a_n=\frac{\partial}{\partial b_{-n}},\qquad
       b_n=-\frac{\partial}{\partial a_{-n}}\quad(n>0).
\end{equation}
With the usual creation and annihilation parts of $a(\xi)$,
\[
 \mathsf Y(a_{-1}f,\xi)
       =a(\xi)_-f(b(\xi))+f(b(\xi))a(\xi)_+.
\]
The field of $b_{-1}h$ is $h(b(\xi))\mathsf T b(\xi)$.
Its two factors have regular OPE.

Put $U=pf+qh$ and $b_-(\xi)=b_0+q\xi+r\xi^2+\cdots$.
Only the annihilator differentiating $p$ acts inside
$f(b(\xi))U$. Therefore
\begin{equation}\label{ch:eq:annihilation-identities}
 \begin{aligned}
 f(b(\xi))U&=f(b_-(\xi))U
                          -\xi^{-1}f'(b_-(\xi))f(b_0),\\
 a(\xi)_+U&=\xi^{-1}(pf'+qh')+\xi^{-2}h.
 \end{aligned}
\end{equation}
Expanding the displayed functions through their first Taylor
terms gives the entire singular part
\begin{equation}\label{ch:eq:self-OPE}
       A(\xi)A(0)\sim
       \frac{C(b_0)}{\xi^2}+
       \frac{q\,C'(b_0)}{2\xi},
       \qquad C=2fh-(f')^2.
\end{equation}
In particular, the simple-pole coefficient is
$q(fh'+f'h-f'f'')$.
The annihilation operators in
\eqref{ch:eq:annihilation-identities} give at most a double pole.
Equivalently, the weight grading permits no coefficient of
negative weight from two weight-one states.
Thus no higher Taylor term can contribute another singularity.
The value of $h$ in \eqref{ch:eq:local-lift} makes $C=0$.

On $B=g(b_0)\mathbf1$, only $a_0$ acts nontrivially among the
nonnegative $a$-modes. Hence
\[
                         A(\xi)B(0)\sim\frac{fg'}\xi=\frac1\xi.
\]
The $b$-fields have regular mutual products, so $B(\xi)B(0)$
is regular. All coefficient functions depend only on $b$.
The spectator pairs consequently have regular OPE with the
changed pair and retain their own mixed OPEs.
\end{proof}

\begin{example}\label{ch:ex:omitted-quantum-term}
For $g(b)=e^b$, the correct coefficients are $f=e^{-b}$ and
$h=e^{-b}/2$.
Omit just the second summand of the state in
\eqref{ch:eq:local-lift}, and write $A^0=a_{-1}e^{-b_0}\mathbf1$.
The mixed product still has its prescribed pole, but
\begin{equation}\label{ch:eq:omitted-self-OPE}
 \begin{aligned}
 A^0(\xi)B(0)&\sim\frac1\xi,\\
 A^0(\xi)A^0(0)&\sim
       -\frac{e^{-2b_0}}{\xi^2}
       +\frac{e^{-2b_0}b_{-1}}\xi.
 \end{aligned}
\end{equation}
This follows from \eqref{ch:eq:self-OPE} with $h=0$.
The double-pole coefficient is nowhere zero on the chart.
Preserving the mixed OPE alone therefore does not define the
coordinate-change homomorphism.
\end{example}

\subsection{Vertex normal ordering, composition and inverses}

The mode order in \eqref{ch:eq:local-lift} can be converted
to a coefficient-first vertex product.
For states $s,t$, write $s_{(-1)}t$ for the constant coefficient
of $\mathsf Y(s,\xi)t$, and use $:f a:$ for that vertex product.
Equation~\eqref{ch:eq:Fock-action} gives
\begin{equation}\label{ch:eq:ordering-conversion}
       f_{(-1)}a_{-1}\mathbf1
         =a_{-1}f(b_0)\mathbf1-f''(b_0)b_{-1}\mathbf1.
\end{equation}
Indeed, the constant term from the annihilation contribution
$-\xi^{-1}f'(b_-(\xi))$ is $-f''q$.
Thus the same lifted field is
\begin{equation}\label{ch:eq:coefficient-first}
                 A(\xi)=:f a:(\xi)
                       +(h+f'')(b(\xi))\mathsf T b(\xi).
\end{equation}
For the exponential change, this coefficient is $3e^{-b}/2$.
Equation~\eqref{ch:eq:coefficient-first} uses the specified
vertex product. It must not be confused with a different
ordering of all fundamental creation and annihilation operators.
The omitted term in Example~\ref{ch:ex:omitted-quantum-term}
was the summand in the original mode expression.

\begin{proposition}\label{ch:prop:composition}
The lifts \eqref{ch:eq:local-lift} compose with holomorphic
one-coordinate changes. The identity has $f=1,h=0$.
On the chosen exponential chart, the exponential and its
logarithm branch give inverse vertex algebra isomorphisms in
both orders. These maps preserve the vacuum and translation.
\end{proposition}

\begin{proof}
For a second coordinate change $k$, let
\[
 f_k=\frac1{k'},\qquad h_k=\frac{(f_k')^2}{2f_k},\qquad
 F=f_k(g(b)),\qquad H=h_k(g(b)).
\]
The constant mode of the first lifted field on $F$ is
\[
             A_{(-1)}F
                 =a_{-1}fF+b_{-1}(hF+f'F').
\]
The extra term $f'F'$ comes from the constant coefficient of
$f(b_-(\xi))F'/\xi$.
The transformed $b_{-1}$ term contributes $b_{-1}g'H$.
The composite coefficients are consequently
\begin{equation}\label{ch:eq:composition}
          f_{\mathrm{comp}}=fF,\qquad
          h_{\mathrm{comp}}=hF+f'F'+g'H.
\end{equation}
Since $g'=1/f$ and $F'=f_k'(g)/f$, the last summand is
$f(F')^2/(2F)$.
Expanding a square gives
\[
        h_{\mathrm{comp}}
             =\frac{(f'F+fF')^2}{2fF}
             =\frac{(f_{\mathrm{comp}}')^2}{2f_{\mathrm{comp}}},
        \qquad f_{\mathrm{comp}}=\frac1{(k\circ g)'}.
\]
Thus the composite has the same canonical formula.

For $g=\exp$ and $k=\log$ on its chosen image,
$F=e^b$ and $H=1/(2e^b)$.
The three terms in $h_{\mathrm{comp}}$ are
\[
                             \tfrac12,\quad -1,\quad \tfrac12.
\]
They sum to zero, and $f_{\mathrm{comp}}=1$.
In the reverse order, with nonzero coordinate $v_{\mathrm{fib}}$,
the three terms are
\[
 \frac1{2v_{\mathrm{fib}}^2},\qquad
 -\frac1{v_{\mathrm{fib}}^2},\qquad
 \frac1{2v_{\mathrm{fib}}^2}.
\]
They again sum to zero, with first coefficient one.
For the identity coordinate map, \eqref{ch:eq:local-lift}
itself gives $f=1,h=0$.

The full generator OPEs in Proposition~\ref{ch:prop:local-OPE},
together with holomorphic functional calculus, determine the
local free vertex algebra homomorphisms.
The two composite calculations identify their inverses.
The constant function remains the vacuum, and the homomorphisms
commute with $\mathsf T$, in particular sending
$\mathsf T b$ to $\mathsf T g(b)$.
Both spectator pairs are fixed.
\end{proof}

\subsection{The local Virasoro law}

For one pair the standard Virasoro state is
$L=a_{-1}b_{-1}\mathbf1=pq\mathbf1$.
Under \eqref{ch:eq:local-lift}, its image is
$\widetilde L=A_{(-1)}(\mathsf T B)$.
The two-jet expansion needed for this constant coefficient is
\[
 f(b_-(\xi))
   =f+f'q\xi+\left(f'r+\tfrac12f''q^2\right)\xi^2+\cdots.
\]
Using the same mode action as above gives
\begin{equation}\label{ch:eq:Virasoro-mode-calculation}
 \widetilde L
   =pqfg'+rf'g'
            +q^2\left(f'g''+\tfrac12f''g'+hg'\right).
\end{equation}
Substitute $f=1/g'$ and $h=(f')^2/(2f)$.
The result is
\begin{equation}\label{ch:eq:local-Virasoro}
       \widetilde L
         =L-(\log g')'r-\tfrac12(\log g')''q^2
         =L-\tfrac12\mathsf T^2\log g'.
\end{equation}
Here $\mathsf T^2F(b_0)=2F'r+F''q^2$.
For $g=\exp$, equation~\eqref{ch:eq:local-Virasoro} reads
$\widetilde L=L-b_{-2}\mathbf1$.
Adding the two unchanged spectator terms proves the same law
for this three-dimensional chart, with its fixed sign and factor.

This local construction checks the quantum correction and its
composition on an actual logarithmic chart.
It gives no formula for arbitrary three-dimensional Jacobians
or for compatible lifts on all overlaps of $Y$.
The global obstruction is the nonzero class
\eqref{ch:eq:nonzero-gerbe}; local quantum corrections cannot
trivialize it. The curve calculation also gives the exact
nonnegative zero cone \eqref{tc:eq:exact-positive-cone} and
the empty pluriclosed Hermitian cone of
Theorem~\ref{tc:thm:no-SKT}.
The standard Virasoro obstruction \eqref{ch:eq:three-kappa}
is logically separate. There is already no global graded
bosonic object to which a standard Virasoro field could be attached.
These conclusions retain the fixed nonzero parameter and selected
filling. The enlarged construction with fermionic fields lies
outside the bosonic gerbe considered here.
Parameter degeneration requires a relative construction.
Comparisons with a field theory on the target or with the finite
Heisenberg representation require their own constructions.

The classical theory in Section~\ref{hp:sec} begins with a different
object: the resolved Hamiltonian observables of a spacetime theory
on the threefold. Its smooth-fiber pushforward can be computed
without a global bosonic target-space CDO sheaf. A relation between
these two constructions would require an additional comparison
of their fields, observables and quantum gluing.
